% GENERATED from QPf5-display.md by md2tex.py. Do not hand-edit: sync is
% one-way and the next run overwrites this file.

\section{The contract, adopted and not forked}
\textbf{The plugin's full tree}: one \texttt{display/} per page, one folder per unit, every unit the same shape.

\begin{verbatim}
  <page>/
    <page>.md                          the page, on stage
    display/                            the plugin . discovery never enters
      <PageID>-Display1-<slug>/        ONE rendered unit
        README.md      claim . kind . caption-job . fragility . status
        intake/         caller-owned, provenance-bound values
          manifest.yaml
          inputs/      small approved summary CSV or JSON extracts only
        recipe/         renderer-owned script, spec, receipts
        float.tex      caption + label + asset reference
        preview.tex    standalone wrapper that inputs float.tex
        preview.pdf     compiled look, what the tab frames
        assets/         the WINNING render
        candidates/    pre-decision renders
        versions/      superseded variants, kept for history
      <PageID>-Display2-<slug>/        the next unit, same shape
\end{verbatim}
The layout above is \texttt{display-unit-output-contract.md}, quoted rather than redesigned. \texttt{intake/} and \texttt{recipe/} are primary material a person rules; \texttt{assets/} and \texttt{preview.pdf} are derived and a rebuild may overwrite them. The one page-side change is the address: the unit root is the page's own \texttt{display/} folder, so the unit travels with the page it illustrates.


\section{The surface: one card per unit in the right pane}
\textbf{The tab}: what a reader sees when they open 🖼 Display on a page.

\begin{verbatim}
  right pane .  Display
    +-----------------------------------------+
    |  QPf5-Display1-drift-bands            |
    |    preview.pdf, framed (body.py ccpdf)  |
    |    README status .  rebuild            |
    |    accepted: [ ] . shown, NOT tickable    |
    +-----------------------------------------+
    one card per unit . empty display/ = empty state, not an error
\end{verbatim}
The tab is the Slides sandwich with a different filling: a drawer plugin posts to a \texttt{live/} endpoint, the endpoint lists the page's units, and the pane frames each \texttt{preview.pdf} through the PDF object path \texttt{src/body.py} already renders. 🔄 rebuild runs the unit's recipe and refreshes the derived half; it never touches \texttt{intake/} and never touches an acceptance row. A unit is EVIDENCE the content cites, and the citation lives IN the sentence: the unit's id named in prose chips as the evidence card in place, carrying its acceptance state and linking its block in the strip. The \texttt{> Display:} lane beside \texttt{> Citation:} and \texttt{> Value:} stays the FILING surface: a machine appending evidence writes a lane and never edits prose, and a binding no sentence carries naturally lands there. Either surface naming a [ ] unit binds a pending render, not accepted evidence. The projections inherit the citation (JL 260816): the latex export embeds a cited unit as a real float after the citing paragraph, built from the winning asset and the unit's own caption, and the word export embeds the rasterized figure with an inline figure number and a 🖼 Display comment on the sentence.


\section{The writers, and the row no machine may tick}
\textbf{Who writes what}: three hands on one unit, each confined to its half.

\begin{verbatim}
   renderer skills   haipipe-display-table . -figure . -diagram . ...
                       write recipe/ . assets/ . preview.pdf
   a person          rules intake/ and ticks accepted:
   the pane          shows everything . writes NOTHING
\end{verbatim}
The renderers are the display family's existing skills, dispatched the way \texttt{task-for-display} already dispatches them; this page adds no renderer. Acceptance follows the model the slide plugin already borrowed: a [x] means a person looked at a specific render and said yes, so a machine that ticks one has forged a judgment, not saved time.

\textbf{The four renderers}: which skill draws which kind, and what each one's recipe is.

\begin{verbatim}
   data-driven . needs a ruled intake/ of numbers
      -table         aggregated CSV/JSON > booktabs LaTeX . recipe = the spec
      -figure        results > plot . recipe = the python script
   concept-driven . the intake is a spec, not numbers
      -diagram       FigureSpec JSON > deterministic editable SVG
      tikz           .tikz.tex recipe > vector PDF . TeX-native: the paper's
                       fonts, math labels, and float.tex may input it directly
      -illustration  image-gen prompt > concept figure . recipe = prompt + receipts
   -poster . -slides  composites, assembling units the renderers above made
\end{verbatim}
A data-driven renderer may not recompute from raw data; it reads the approved extract in \texttt{intake/} and nothing else. A concept-driven renderer carries no numbers, so its intake is the spec it draws. The rebuild differs the same way: a figure rebuild reruns its python script, a diagram rebuild is deterministic from its spec, and an illustration rebuild replays a receipted prompt.


\section{How a unit is generated}
\textbf{The pipeline}: five steps from approved numbers to an accepted render, and who moves each one.

\begin{verbatim}
  (1)  INTAKE      the caller fills intake/: manifest.yaml + small
                     approved extracts, every value provenance-bound
  (2)  RENDER     the matching renderer skill runs, dispatched through
                     task-for-display; it fills recipe/ + candidates/
  (3)  PICK       a person picks the winner; it moves to assets/,
                     the losers stay in candidates/ or versions/
  (4)  BUILD      preview.tex compiles to preview.pdf; the adapter's
                     --check reports a stale unit instead of hiding it
  (5)  ACCEPT     a person reads the render and ticks accepted:
\end{verbatim}
The pipeline never starts at (2): a renderer given no ruled \texttt{intake/} has nothing provenance-bound to draw, which is what \texttt{display-intake-contract.md} exists to refuse. Steps (1) (3) (5) are a person's and steps (2) (4) are machinery, the split of hands QPf5-Display1 draws; 🔄 rebuild on the tab is step (4) alone, rerun. A changed \texttt{intake/} flows forward, (2) to (5), and the acceptance falls back to [ ] because the bytes a person said yes to are gone.

\begin{figure}[t]
\centering
\includegraphics[width=.85\linewidth]{../display/QPf5-Display1-pipeline-tikz/assets/figure.pdf}
\caption{The display pipeline: five steps from approved intake to an accepted
  render. Green steps belong to a person, blue steps to machinery, and a changed
  intake returns the acceptance to open.}
\label{fig:qpf5-pipeline}
\end{figure}

\textbf{Where (1)'s numbers come from}: the evidence chain behind an intake, id-bound at every hop.

\begin{verbatim}
   the page needs a number   > raised as a PROBE . the page's own probe/ plugin (QPf9)
   probe orchestrator        > sweeps the bank . reuse | enrich | fresh
   haipipe-task runs         > answer lands in <task-folder>/QA/<n>-<slug>.md
                                 + the run's artifact, e.g. counts.csv . CANONICAL
   intake/manifest.yaml      > cites that holder BY ID . carries only a
                                 small approved extract in inputs/
\end{verbatim}
\texttt{intake/} is not a second data store: the task output stays canonical, and the manifest points back at the exact holder and run that produced it. A unit names its evidence by id, never by path, and never types a number; QPf5-Display2 draws this ask-once-cite-twice move. The slide plugin already walks this chain on \texttt{QPf3}: its deck declares \texttt{needs: QA-probe/QBt5-for-value/1-artifact-paths}, the resolver turns the id into the artifact's path, and the template carries no digits, so a hand-typed value has nowhere to live. A1's adapter is what makes the same id resolution work for a unit under a page.

\begin{figure}[t]
\centering
\includegraphics[width=.85\linewidth]{../display/QPf5-Display2-small-paper-tikz/assets/figure.pdf}
\caption{A page as a small paper: the probe plugin asks a question once, the
  bank answers behind the wall, and both the display intake and the page's own
  sentence cite the same card by id.}
\label{fig:qpf5-small-paper}
\end{figure}

